\documentclass[12pt, a4paper]{article}

% ── Packages ──────────────────────────────────────────────────
\usepackage[T1]{fontenc}
\usepackage[utf8]{inputenc}
\usepackage{lmodern}
\usepackage{microtype}
\usepackage[top=2.5cm, bottom=2.5cm, left=3cm, right=3cm]{geometry}
\usepackage{amsmath, amssymb}
\usepackage{booktabs}
\usepackage{array}
\usepackage{parskip}
\usepackage{titlesec}
\usepackage{enumitem}
\usepackage[hidelinks]{hyperref}

\titleformat{\section}{\large\bfseries}{\thesection.}{0.5em}{}[\titlerule]
\titleformat{\subsection}{\normalsize\bfseries}{\thesubsection.}{0.5em}{}

\newcolumntype{L}[1]{>{\raggedright\arraybackslash}p{#1}}

% ──────────────────────────────────────────────────────────────
\begin{document}

\begin{center}
    {\LARGE\bfseries Glossary of Terms and Acronyms}\\[0.5em]
    {\large Transfer Learning for Freshness Classification}\\[0.3em]
    {\normalsize Digital Image Processing --- 2026}
\end{center}

\vspace{1em}
\hrule
\vspace{1.5em}

% ─────────────────────────────────────────────────────────────
\section{Acronyms and Abbreviations}
% ─────────────────────────────────────────────────────────────

\begin{table}[h]
\begin{tabular}{L{2.5cm} L{11cm}}
\toprule
\textbf{Acronym} & \textbf{Definition} \\
\midrule

AUC-ROC & \textit{Area Under the Receiver Operating Characteristic Curve.}
          A metric ranging from 0.5 (random classifier) to 1.0 (perfect
          classifier) that measures the model's ability to separate classes
          across all decision thresholds, independent of a fixed threshold. \\[4pt]

CNN     & \textit{Convolutional Neural Network.}
          A class of deep learning model that applies learned filters to
          detect local patterns (edges, textures, shapes) in images.
          The foundational architecture for computer vision tasks. \\[4pt]

C1--C4  & \textit{Training Conditions 1 through 4.}
          An ablation study over backbone fine-tuning and head architecture.
          See Section~\ref{sec:conditions}. \\[4pt]

DIP     & \textit{Digital Image Processing.}
          The field of study concerned with the manipulation and analysis
          of digital images using computational methods. \\[4pt]

EB0     & \textit{EfficientNet-B0.}
          A compact convolutional backbone (5.3M parameters) achieving
          77.1\% ImageNet top-1 accuracy through compound scaling. \\[4pt]

EB2     & \textit{EfficientNet-B2.}
          A larger EfficientNet variant (9.1M parameters) achieving
          80.1\% ImageNet top-1 accuracy. \\[4pt]

F1      & \textit{F1-Score (macro-averaged).}
          The harmonic mean of precision and recall, averaged equally
          across all classes. The primary evaluation metric in this study. \\[4pt]

FN      & \textit{False Negative.}
          A rotten sample incorrectly classified as fresh.
          The most dangerous error in food safety applications. \\[4pt]

FP      & \textit{False Positive.}
          A fresh sample incorrectly classified as rotten.
          Results in unnecessary food waste. \\[4pt]

GPU     & \textit{Graphics Processing Unit.}
          Hardware accelerator used to parallelize the matrix operations
          required for deep learning training and inference. \\[4pt]

ImageNet & \textit{Large Scale Visual Recognition Challenge dataset.}
           A benchmark dataset of 1.2 million images across 1,000 categories,
           used to pre-train backbone architectures for transfer learning. \\[4pt]

KFQ     & \textit{Kaggle Fruits Quality.} Dataset code. 359 images, 12 mixed
          fruit categories, binary fresh/rotten labels. \\[4pt]

KFR     & \textit{Kaggle Fruits Fresh/Rotten.} Dataset code. 13,599 images,
          apple/banana/orange, binary fresh/rotten labels. \\[4pt]

KFS     & \textit{Kaggle Fresh \& Stale.} Dataset code. 27,317 images,
          9 fruit and vegetable categories, binary fresh/rotten labels. \\[4pt]

LR      & \textit{Learning Rate.}
          The step size used by the optimizer when updating model weights.
          Controls the speed and stability of training convergence. \\[4pt]

MCC     & \textit{Matthews Correlation Coefficient.}
          A balanced metric ranging from $-1$ (inverse prediction)
          through $0$ (random) to $+1$ (perfect prediction).
          More robust than F1-score for imbalanced datasets. \\[4pt]

MFR     & \textit{Mendeley Fruits Classification.} Dataset code. 1,655 images,
          peach/pomegranate/strawberry, binary fresh/rotten labels. \\[4pt]

MFV     & \textit{Mendeley FruitVision.} Dataset code. 10,154 images,
          5 fruit categories, three-class fresh/formalin/rotten labels. \\[4pt]

MLM     & \textit{Mendeley Lemon Varieties.} Dataset code. 1,956 images,
          7 lemon varieties, binary fresh/rotten labels. \\[4pt]

MN3     & \textit{MobileNetV3-Small.}
          A lightweight backbone (2.5M parameters) optimized for
          mobile and edge deployment scenarios. \\[4pt]

OWN     & \textit{Own Dataset.} 121 real strawberry photos taken with a
          mobile phone camera under natural lighting conditions.
          Used exclusively for cross-domain evaluation. \\[4pt]

R18     & \textit{ResNet-18.}
          A residual network backbone (11.2M parameters) achieving
          69.8\% ImageNet top-1 accuracy. Standard transfer learning baseline. \\[4pt]

ReLU    & \textit{Rectified Linear Unit.}
          Non-linear activation function defined as $f(x) = \max(0, x)$.
          Applied between layers of the projection head. \\[4pt]

ROC     & \textit{Receiver Operating Characteristic curve.}
          A plot of true positive rate vs.\ false positive rate at all
          classification thresholds. See AUC-ROC. \\[4pt]

ST      & \textit{State mode.}
          Label mode where classes correspond to freshness states only:
          fresh, rotten, or formalin. \\[4pt]

TN      & \textit{True Negative.}
          A rotten sample correctly classified as rotten. \\[4pt]

TP      & \textit{True Positive.}
          A fresh sample correctly classified as fresh. \\[4pt]

\bottomrule
\end{tabular}
\end{table}

% ─────────────────────────────────────────────────────────────
\section{Key Concepts}
% ─────────────────────────────────────────────────────────────

\subsection{Transfer Learning}

Transfer learning is a machine learning technique in which a model
pre-trained on a large dataset (source domain) is adapted to perform
a different but related task on a smaller dataset (target domain).
In this study, all backbone models are pre-trained on ImageNet and
then fine-tuned for fruit and vegetable freshness classification.
The core assumption is that low-level features (edges, textures, colors)
and mid-level features (shapes, patterns) learned from natural images
transfer well to the food quality domain.

\subsection{Backbone Architecture}

A backbone is the main feature extraction component of a deep learning
model for image classification.
It consists of a series of convolutional layers that progressively
transform a raw image into a compact feature vector.
In transfer learning, the backbone is initialized with pre-trained weights
and may be partially or fully frozen during training on the target task.
The output of the backbone --- typically obtained after global average
pooling --- is a fixed-length vector (e.g.\ 512 dimensions for ResNet-18)
that summarizes the image content.

\subsection{Fine-tuning}

Fine-tuning refers to the process of continuing to train a pre-trained
model on new data, allowing some or all of its weights to be updated.
\textit{Partial fine-tuning} (as in conditions C2 and C4) unfreezes
only the last few layers of the backbone, which capture the most
task-specific high-level features.
\textit{Full fine-tuning} updates all weights, which risks
catastrophic forgetting unless the learning rate is carefully controlled.

\subsection{Catastrophic Forgetting}

A phenomenon in neural networks where learning new information
causes the model to ``forget'' previously acquired knowledge.
In transfer learning, aggressive fine-tuning of a pre-trained backbone
destroys the general-purpose representations learned during pre-training,
eliminating the benefit of the pre-trained initialization.
This is mitigated in condition C4 by using a backbone learning rate
100 times smaller than the head learning rate
($\text{lr}_\text{backbone} = 10^{-5}$ vs $\text{lr} = 10^{-3}$).

\subsection{Projection Head}

A non-linear transformation applied between the backbone feature vector
and the final linear classifier.
In conditions C3 and C4, the projection head consists of two fully
connected layers with ReLU activation:
\[
  \mathbf{f} \in \mathbb{R}^{d_\text{out}}
  \xrightarrow{\text{FC}(256)} \xrightarrow{\text{ReLU}}
  \xrightarrow{\text{FC}(128)} \xrightarrow{\text{ReLU}}
  \mathbf{h} \in \mathbb{R}^{128}
\]
The projection head provides additional non-linear capacity to form
more complex decision boundaries in the feature space, without requiring
backbone adaptation.

\subsection{Domain Shift}

Domain shift refers to the statistical difference between the distribution
of training data and the distribution of data encountered at test time.
In this study, all backbone models are trained on internet-sourced
images (typically studio photographs with controlled lighting and
white backgrounds) and evaluated on real-world mobile phone photos
taken under natural indoor lighting conditions.
The resulting performance degradation --- measured as
$\text{Drop} = F1_\text{train} - F1_\text{eval}$ --- quantifies
the magnitude of domain shift.

\subsection{Macro-Averaged Metrics}

When a classification task involves multiple classes, metrics such as
F1-score, precision, and recall can be computed per class and then
averaged.
\textit{Macro averaging} computes the metric for each class independently
and takes the unweighted mean:
\[
  \overline{F1} = \frac{1}{C} \sum_{c=1}^{C} F1_c,
  \quad
  F1_c = \frac{2\,TP_c}{2\,TP_c + FP_c + FN_c}
\]
This gives equal weight to every class regardless of its sample count,
which is important when minority classes (such as formalin) are
particularly critical from a safety perspective.

\subsection{Confusion Matrix}

A square matrix of size $C \times C$ where entry $(i, j)$ counts
the number of samples of true class $i$ predicted as class $j$.
The diagonal entries represent correct predictions; off-diagonal entries
represent misclassifications.
In this study, confusion matrices are \textit{row-normalized} so that
each entry shows the fraction of samples in that true class, making
error patterns directly comparable across classes of different sizes.

\subsection{Cosine Annealing}

A learning rate scheduling strategy in which the learning rate follows
a cosine curve from its initial value down to zero over the training
period:
\[
  \text{lr}(t) = \frac{\text{lr}_\text{max}}{2}
  \left(1 + \cos\!\left(\frac{\pi\, t}{T}\right)\right)
\]
where $t$ is the current epoch and $T$ is the total number of epochs.
This allows large gradient steps early in training and increasingly
fine adjustments as the model approaches convergence.

\subsection{Stratified Split}

A data splitting strategy that ensures each subset (train and validation)
contains the same proportion of samples from each class as the full
dataset.
In this study, an 80/20 stratified split is applied with a fixed
random seed (42) to ensure reproducibility.
This prevents scenarios where, by chance, one class is over- or
under-represented in the validation set.

\subsection{Formalin Treatment}

Formalin (formaldehyde solution) is used in some markets to preserve
the appearance of fruits and vegetables, making them visually similar
to fresh produce despite being unsafe for consumption.
The FruitVision dataset includes formalin-treated samples as a
third class, making it the most challenging benchmark in this study.
The difficulty of distinguishing formalin-treated from fresh samples
in RGB images is directly analogous to the challenge of detecting
subtle spoilage in fermented Korean foods.

% ─────────────────────────────────────────────────────────────
\section{Training Condition Reference}
\label{sec:conditions}
% ─────────────────────────────────────────────────────────────

\begin{table}[h]
\centering
\begin{tabular}{L{1.5cm} L{2.5cm} L{3cm} L{7cm}}
\toprule
\textbf{ID} & \textbf{Backbone} & \textbf{Head} & \textbf{Full description} \\
\midrule
C1 & Fully frozen & Linear only &
  All backbone weights locked. Single linear layer $\mathbf{W} \in
  \mathbb{R}^{d_\text{out} \times C}$ trained from scratch.
  Evaluates raw ImageNet feature quality. \\[6pt]

C2 & Last block free & Linear only &
  \texttt{layer4} (final residual block) unfrozen.
  All earlier backbone layers frozen.
  Linear classifier only, no projection head.
  Allows high-level feature adaptation. \\[6pt]

C3 & Fully frozen & Proj.\ + Linear &
  All backbone weights locked.
  Two-layer projection head:
  $d_\text{out} \to 256 \to 128$, each followed by ReLU,
  then linear classifier.
  Non-linear capacity without backbone adaptation. \\[6pt]

C4 & Last block free & Proj.\ + Linear &
  \texttt{layer4} unfrozen with $\text{lr}_\text{backbone} = 10^{-5}$.
  Full projection head with $\text{lr} = 10^{-3}$.
  Maximum adaptation capacity with catastrophic forgetting prevention. \\[6pt]

\bottomrule
\end{tabular}
\end{table}

% ─────────────────────────────────────────────────────────────
\section{Dataset Reference}
% ─────────────────────────────────────────────────────────────

\begin{table}[h]
\centering
\begin{tabular}{L{1.2cm} L{3.5cm} L{1.5cm} L{2cm} L{4cm}}
\toprule
\textbf{ID} & \textbf{Full name} & \textbf{Images} & \textbf{Classes} & \textbf{Notes} \\
\midrule
KFQ & Kaggle Fruits Quality       &    359 & fresh / rotten & 12 mixed fruit types \\
MFR & Mendeley Fruits              &  1,655 & fresh / rotten & Peach, pomegranate, strawberry \\
MLM & Mendeley Lemon Varieties     &  1,956 & fresh / rotten & 7 lemon varieties \\
MFV & Mendeley FruitVision         & 10,154 & fresh / formalin / rotten & Most challenging; 3 classes \\
KFR & Kaggle Fruits Fresh/Rotten   & 13,599 & fresh / rotten & Apple, banana, orange \\
KFS & Kaggle Fresh \& Stale        & 27,317 & fresh / rotten & 9 fruit/vegetable categories \\
OWN & Own Dataset (strawberry)     &    121 & fresh / rotten & Real phone photos; eval only \\
\midrule
    & \textbf{Total (public)}      & \textbf{55,040} & & \\
\bottomrule
\end{tabular}
\end{table}

\end{document}